\documentclass[onecolumn, draftclsnofoot,10pt, compsoc]{IEEEtran}
\usepackage{graphicx}
\usepackage{url}
\usepackage{setspace}

\usepackage{geometry}
\geometry{textheight=9.5in, textwidth=7in}

% 1. Fill in these details
\def \CapstoneTeamName{		The Cleverly Named Team}
\def \CapstoneTeamNumber{		63}
\def \GroupMemberOne{			Joshua Lioy}
\def \GroupMemberTwo{			Corynna Park}
\def \GroupMemberThree{			Jackson Wells}
\def \CapstoneProjectName{		Molecules in 3D?! And in color!? That I can hold in my hand? No way!!!}
\def \CapstoneSponsorCompany{		Oregon State University}
\def \CapstoneSponsorPerson{		Victor Hsu}

% 2. Uncomment the appropriate line below so that the document type works
\def \DocType{			Problem Statement
				%Requirements Document
				%Technology Review
				%Design Document
				%Progress Report
				}
			
\newcommand{\NameSigPair}[1]{\par
\makebox[2.75in][r]{#1} \hfil 	\makebox[3.25in]{\makebox[2.25in]{\hrulefill} \hfill		\makebox[.75in]{\hrulefill}}
\par\vspace{-12pt} \textit{\tiny\noindent
\makebox[2.75in]{} \hfil		\makebox[3.25in]{\makebox[2.25in][r]{Signature} \hfill	\makebox[.75in][r]{Date}}}}
% 3. If the document is not to be signed, uncomment the RENEWcommand below
\renewcommand{\NameSigPair}[1]{#1}

%%%%%%%%%%%%%%%%%%%%%%%%%%%%%%%%%%%%%%%
\begin{document}
\begin{titlepage}
    \pagenumbering{gobble}
    \begin{singlespace}
    	%\includegraphics[height=4cm]{./coe_v_spot1}
        \hfill 
        % 4. If you have a logo, use this includegraphics command to put it on the coversheet.
        %\includegraphics[height=4cm]{CompanyLogo}   
        \par\vspace{.2in}
        \centering
        \scshape{
            \huge CS Capstone \DocType \par
            {\large\today}\par
            \vspace{.5in}
            \textbf{\Huge\CapstoneProjectName}\par
            \vfill
            {\large Prepared for}\par
            \Huge \CapstoneSponsorCompany\par
            \vspace{5pt}
            {\Large\NameSigPair{\CapstoneSponsorPerson}\par}
            {\large Prepared by }\par
            Group\CapstoneTeamNumber\par
            % 5. comment out the line below this one if you do not wish to name your team
            %\CapstoneTeamName\par 
            \vspace{5pt}
            {\Large
                \NameSigPair{\GroupMemberOne}\par
                \NameSigPair{\GroupMemberTwo}\par
                \NameSigPair{\GroupMemberThree}\par
            }
            \vspace{20pt}
        }
        \begin{abstract}
        % 6. Fill in your abstract
	    Our project aims to produce a robust workflow for 3D printing multi-colored models. 
	    Currently, the capability to produce single-colored 3D models already exists. 
	    We look to expand upon this by adding a multi-color functionality.
	    This will enable enhanced learning opportunities and better absorption of concepts, specifically those relating to the field of biochemistry. 
	    While this project originated as one for the OSU Biochemistry department, we hope that our findings about how to perform the process of multi-color 3D printing will benefit those looking to 3D pring multi-color models in general, not just those relating to the field of Biochemistry.

	    To complete this project, we will need to become familiar with the existing workflow of how to print single-colored models. 
	    From there, we will utilize existing free and open-source softwares to deduce the best way to create multi-color 3D models.
	    At the completion of this project, we will have produced a robust workflow as well as a journal article describing the process we created to utilize 3D multi-colored printing. 

        	%This document is written using one sentence per line.
        	%This allows you to have sensible diffs when you use \LaTeX with version control, as well as giving a quick visual test to see if sentences are too short/long.
        	%If you have questions, ``The Not So Short Guide to LaTeX'' is a great resource (\url{https://tobi.oetiker.ch/lshort/lshort.pdf})
        \end{abstract}     
    \end{singlespace}
\end{titlepage}
\newpage
\pagenumbering{arabic}
\tableofcontents
% 7. uncomment this (if applicable). Consider adding a page break.
%\listoffigures
%\listoftables
\clearpage

%\newpage

% 8. now you write!
\section{Defining the Problem}
A technology that once used to be available only by the highly affluent or highly qualified individuals, 3D printing is becoming more and more accessible to the general public, and almost anyone can learn how to utilize a 3D printer. 
Printing in 3D opens the door to enhanced learning opportunities, especially for the students that learn best through kinesthectic means.
Kinesthetic or tactile learning is learning in which concepts are understood through physical activities such as touching, feeling, and holding. 
For the problem defined by our client, the field in focus is that of biochemistry. 
Courses in biochemistry stand to benefit greatly from the ability to incorporate 3D printing in course material. 
Specifically, being able to 3D print models of various molecular structures, for example, DNA, different atoms such as carbon, nitrogen, and oxygen, and proteins can greatly aid the understanding of students regarding important concepts in the field of biochemistry.

Currently, 3D printing molecular structures in a single color is doable. 
There already exists many open source or free programs which take molecular structures and create 3D printable STL files.
Furthermore, there are several journal articles online which provide guidance and a robust workflow on how to successfully peform the process of 3D printing.  
However, as to be expected, there are ways to improve. 
Our client Victor Hsu, representing the needs of Oregon State University, believes that while 3D printing single-color models has already proven to be of significant worth when it comes to student learning, more could be obtained if models were able to be printed in multiple colors. 
Colors can be used to provide more information about biomolecular structure, such as atom types, hydrophobic surfaces, etc. 

Up until now, in order to produce models that convey information with color, the OSU Biochemistry department print 3D models using white filament and then hand-paint individual atoms to indicate atom type. 
Hsu sees multi-color 3D printing as the solution to this tedious and difficult process.
Recently, the OSU Biochemistry department has purchased two 3D printers for the purpose of 3D printing molecular structures. 
They also have purchased a multimaterial printing interface that will allow 3D printing in up to 4 different colors. 
Multi-filament printers are only now becoming available, thus, there exists little to no documentation on how to use them to produce intended results. 
Thus, the problem that stands is the need for a streamlined workflow for 3D printing molecular models with multiple colors, or any multi-colored 3D model.
By solving this problem, we and our client hope to enhance the scientific learning of students as well as help others achieve a greater understanding and capability regarding 3D printing. 

\section{Proposed Solution}
Our proposed solution is to develop a robust workflow for multicolor 3D printing. 
Additionally, to maximize the benefit of our project accomplishments to the public, provide online documentation of the multi-color 3D printing process by writing a journal article describing the process and then submitting it for publication. 
To achieve these goals, we will need to first become well informed about the 3D printing process, first in a single color. 
We will need to explore the various options we have regarding software tools to generate 3D printable STL files.
From our current understanding, when trying to print a model with multiple colors, multiple STL files are generated.
One approach to the problem could be to find a way to link together the multiple STL files to create one multi-color STL file which would then produce a multi-color model.
We will also need to decide the best file type to use to show and produce a multi-colored model; the two dicussed with our client were STL and OBJ. 

\section{Performance Metrics}
Once the project has been completed, we will be able to use the 3D printers purchased by OSU's Biochemistry Department to produce 3D models in more than one color.
Our aim is to be able to print 3D models with anywhere from one to four colors. 
Should we reach our goal of four different colors, we may then choose to explore the possibility of printing up to eight different colors in one 3D model. 
We will have also produced a robust workflow for multi-color 3D printing that we and the client will be able to follow and produce consistent, desirable results. 
After we produce a satisfactory workflow, we will write a journal article describing the multi-color 3D printing process with the intention that if an individual is wondering how to do the same or a similar process, he or she will be able to benefit from our article. 
These are the metrics that will determine if we have reached completion regarding this project. 

\section{Conclusion}
In conclusion, our project stands to benefit several parties, from the educator to the curious. 
Personally, I look forward to being able to help kinesthetic learners learn to their best ability thanks to the advancements in 3D printing. 
With the ability to 3D print models in color, individuals will be able to have a better physical grasp on many concepts, and not just concepts limited to the field of biochemisty.
We are not just creating a workflow for chemists and biologists, or instructors and students.
We are creating a workflow to create a greater desire for 3D printing and the benefits that they can bring. 
This project aims to improve an already impressive and useful technology and make it easier to use by providing a streamlined and proven process.
As such, it will simultaneously improve user experience and achievement regarding 3D printing. 

\end{document}
